\documentclass{situationist-preprint}

\title{Echolocation and Spatial Experience}
\msid{MND-03-b}
\msfamily{MND-03}
\msdiscipline{Philosophy of Mind}
\mstype{Companion manuscript}
\mskeywords{echolocation; spatial experience; bats; human echolocation; mirroring; sensory substitution}

\begin{document}
\maketitle

\begin{abstract}
MND-03 argues that what it is like to be an echolocating bat is like having a sonar-structured spatial field. A natural objection is that echolocation is a form of hearing, and that hearing in humans presents sounds, not surfaces; the bat's experience should then be auditory in structure rather than spatial. We argue that the structure of an echolocative experience is fixed by what the creature does with the echoes, not by the energy that carries them. We support the claim with the case of human echolocators, whose reports and neural activity show echolocation used as a spatial sense, and show that the case confirms the mirroring principle on which MND-03 relies while leaving the intrinsic residue in place.
\end{abstract}

\section{Hearing sounds and hearing surfaces}

When a human hears a bell, what is presented is a sound, located more or less roughly in a direction, with pitch, loudness and timbre. If the bat's echolocation were hearing in this sense, its experience would present echoes: sounds coming back, varying in delay and spectrum. The bat would then infer the layout of its surroundings from the character of the echoes, and its experience would be auditory, with spatial content added by inference. MND-03 denies this. Lemmas~5.1 and 5.2 of MND-03-a hold that the bat's echolocative experience presents a spatial field of surfaces, and that the field is sonar-structured: range is presented directly, sampled in pulses, and resolved more finely as the bat closes on a target.

The argument for these lemmas is the mirroring principle (MND-03-a, Principle~P2): if a conscious creature uses a modality to guide flexible action, the phenomenal structure of the modality mirrors the structure of the discriminations it uses. The bat uses echo delay to guide pursuit and avoidance as a measure of range, and integrates delay, binaural differences and spectral cues into a representation of a common layout \citep{Griffin1958,Simmons1989}. By the principle, delay is experienced as range, and the integrated values are presented in a common field. The objection under discussion rejects this application, holding that the energy carrying the information, sound, fixes the experiential category.

\begin{proposition}\label{p:carrier}
The structure of a modality's experience is not fixed by the form of energy its receptors transduce.
\end{proposition}

\begin{proof}
By the mirroring principle, the structure is fixed by the discriminations the creature uses and by their integration. Two creatures can transduce the same energy and use different discriminations, or transduce different energies and use the same ones; the principle then assigns them different structures in the first case and the same in the second. The principle would conflict with Proposition~\ref{p:carrier} only if the discriminations used were fixed by the energy transduced, which the following case shows they are not.
\end{proof}

\section{Human echolocators}

Some blind humans echolocate by producing tongue clicks and attending to the echoes. Experts can detect, localize and discriminate objects by size, shape and material, and use the skill to walk, cycle and navigate unfamiliar spaces. \citet{Thaler2011} found that, in two expert echolocators, listening to recordings of their own clicks and echoes activated regions of occipital cortex that are visual in sighted subjects, while activation of auditory cortex did not distinguish echo-bearing from echo-free recordings. The echolocators described what they perceived in spatial terms: objects at locations, with extents and surfaces.

The case bears on Proposition~\ref{p:carrier} in two ways. First, the energy is the same as in ordinary human hearing, but the discriminations used are different: the echolocator uses delay and spectral change as measures of distance, size and texture, and integrates them into a layout through which she moves. By the mirroring principle, her echolocative experience presents a spatial field, and her reports say that it does. Second, the neural evidence fits the prediction that the processing of echoes for spatial guidance is organized like spatial vision rather than like the perception of sounds as such. Neither finding is decisive alone, but together they show that the auditory carrier does not fix an auditory experiential structure. This is the claim needed for Lemma~5.1 of MND-03-a, and the human case gives it support from subjects who can report.

\begin{objection}
Human echolocators do not hear surfaces; they hear echoes and have learned to judge what produced them, as a sommelier judges vineyards. Their spatial reports describe their judgments, not their experience; and a parallel doubt about whether the bat's processing yields a presentation of objects at all is raised by \citet{Akins1993}.
\end{objection}

\begin{reply}
The sommelier's judgment is slow, reflective and detachable from action; the echolocator's use of echoes is fast, continuous and integrated into the control of movement, which is the condition under which the mirroring principle applies. The objection's reading also predicts that the echolocator would be aware first of the echoes and only then of the objects, while expert echolocators report the opposite: the objects are what they notice, and attention to the echoes as sounds is a separate act. The case is closer to the experience of a tactile-visual substitution user who, after training, reports objects out in space rather than patterns on the skin \citep{BachyRita1969}. The bat, whose echolocation is far more refined and has been used for millions of years, is a still clearer case.
\end{reply}

\begin{objection}
If the human echolocator has a sonar-structured spatial field, then we know what it is like to be a bat by asking her, and Nagel's point is refuted entirely. MND-03 says it is not.
\end{objection}

\begin{reply}
The human echolocator shares with the bat the structure of a spatial field presented through echoes, at a coarser grain, and lacks the bat's control of pulse rate, its range resolution and, in some species, its presentation of relative velocity. Even where the structures match, Theorem~5.3 and Proposition~5.4 of MND-03-a hold that what can be shared is the structure, and that the intrinsic character of the qualities that occupy the field cannot be read off from it. Whether the bat's range qualities resemble the echolocator's intrinsically is not settled by their structural match. The human case narrows the residue \citet{Nagel1974} pointed to; it does not remove it.
\end{reply}

\section{Conclusion}

Echolocation is a spatial sense, whatever the energy it uses. The structure of echolocative experience is the structure of the discriminations the creature uses to guide action, which for the bat is a sonar-structured spatial field. Human echolocators confirm that an auditory carrier can support a spatial experience, and they confirm the limits of MND-03 too: structure can be shared and described, intrinsic character cannot.

\begin{thebibliography}{9}\small
\bibitem[Akins(1993)]{Akins1993} Akins, K. (1993). What is it like to be boring and myopic? In B.~Dahlbom (ed.), \emph{Dennett and His Critics}. Blackwell.
\bibitem[Bach-y-Rita et~al.(1969)]{BachyRita1969} Bach-y-Rita, P., Collins, C.~C., Saunders, F.~A., White, B. and Scadden, L. (1969). Vision substitution by tactile image projection. \emph{Nature}, 221, 963--964.
\bibitem[Griffin(1958)]{Griffin1958} Griffin, D.~R. (1958). \emph{Listening in the Dark: The Acoustic Orientation of Bats and Men}. Yale University Press.
\bibitem[Nagel(1974)]{Nagel1974} Nagel, T. (1974). What is it like to be a bat? \emph{The Philosophical Review}, 83(4), 435--450.
\bibitem[Simmons(1989)]{Simmons1989} Simmons, J.~A. (1989). A view of the world through the bat's ear: the formation of acoustic images in echolocation. \emph{Cognition}, 33(1--2), 155--199.
\bibitem[Thaler et~al.(2011)]{Thaler2011} Thaler, L., Arnott, S.~R. and Goodale, M.~A. (2011). Neural correlates of natural human echolocation in early and late blind echolocation experts. \emph{PLoS ONE}, 6(5), e20162.
\end{thebibliography}

\end{document}
